\chapter{Planification de saisie}
\label{ch:planning}

Le chapitre précédent (chapitre~\ref{ch:perception}) a livré une description
géométrique de la scène : pour chaque objet cible, une position 3D dans le
repère base, une boîte englobante, une classe de pose (debout, couché, compact)
et, lorsque l'empreinte le permet, l'orientation de son grand axe. La question
qui s'ouvre maintenant est celle de la décision : \emph{où} et \emph{comment}
poser la pince pour saisir cet objet. C'est l'objet du \textit{grasp planning}.
Ce module est volontairement séparé du contrôle : il décide d'une
\emph{géométrie} de saisie (des poses d'effecteur dans \(SE(3)\)), sans rien
savoir de la cinématique du bras ; c'est le chapitre suivant
(chapitre~\ref{ch:control_eval}) qui traduira ces poses en mouvement
articulaire. Cette frontière nette est ce qui rend chacune des deux étapes
testable isolément.

\section{Définition du problème}
\label{sec:plan_pb}

Étant donné une \texttt{ObjectInstance} fournie par le module de perception, la
planification de saisie doit produire trois poses successives de l'effecteur,
chacune une matrice homogène \(4\times4\) de \(SE(3)\) exprimant
\(T_{\text{base}}^{\text{gripper}}\) (position et orientation de la pince dans le
repère base) :

\begin{description}
  \item[\textit{approach}] pose où la pince est positionnée \emph{au-dessus} de
        l'objet, prête à descendre, gripper ouvert.
  \item[\textit{grasp}] pose où la pince est positionnée \emph{sur} l'objet, au
        moment où le gripper se referme.
  \item[\textit{retract}] pose où la pince est remontée après la saisie, gripper
        fermé maintenant l'objet.
\end{description}

À ces trois poses s'ajoutent deux paramètres scalaires : l'ouverture du gripper
avant l'approche et sa fermeture pour la saisie, exprimés en pourcentage de la
course du gripper. L'ensemble forme une \texttt{GraspPose}
(structure définie dans \texttt{src/planning/grasp.py}).

Ce découpage en \textit{approach / grasp / retract} n'est pas arbitraire : il
sépare la phase de \emph{positionnement} (descendre verticalement sur l'objet
sans le heurter de côté) de la phase de \emph{préhension} (fermer les doigts) et
de la phase de \emph{dégagement} (soulever l'objet une fois saisi). Cette
décomposition en sous-mouvements rectilignes est celle que retient le survey de
Bohg \textit{et al.}~\cite{bohg_data-driven_2014}, et c'est elle qui permet
ensuite au module de contrôle d'enchaîner des trajectoires lisses entre points
de passage connus.

\section{Stratégie retenue : \textit{top-down grasp}}
\label{sec:plan_topdown}

\subsection{Pourquoi le top-down, et pas autre chose}

Bohg \textit{et al.}~\cite{bohg_data-driven_2014} distinguent trois familles de
\textit{grasp planners} : \textit{heuristic} (règles géométriques explicites),
\textit{contact-based} (analyse des points de contact et de la fermeture de
forme), et \textit{data-driven} (synthèse apprise à partir d'exemples). Le choix
d'une approche \textbf{heuristique \textit{top-down}} découle directement de la
question posée et de la plateforme :

\begin{itemize}
  \item Le cahier des charges demande des \emph{stratégies adaptées à la
        géométrie et à l'accessibilité} ; une règle géométrique explicite y
        répond et reste lisible.
  \item La quasi-totalité des objets cibles du TFE sont des objets du quotidien
        posés sur une table (cube, cylindre, mug, rubik's cube), donc
        \emph{accessibles par le dessus}.
  \item Le caractère déterministe et explicable autorise l'analyse d'un échec :
        chaque pose produite peut être inspectée et reliée à une cause.
  \item La méthode fournit une \textbf{base de référence reproductible} pour les
        comparaisons ultérieures avec des méthodes apprises
        (chapitre~\ref{ch:control_eval}).
\end{itemize}

Une raison plus profonde justifie ce choix, et elle n'est pleinement établie
qu'au chapitre de discussion : sur un bras à 5~degrés de liberté actionnés
comme le SO-101, l'approche verticale est la seule à être véritablement
bien conditionnée. Une prise latérale « vraie » (pince horizontale se refermant
autour du flanc vertical d'un objet) est, elle, mal conditionnée sur cette
cinématique. Ce résultat est montré numériquement au
chapitre~\ref{ch:control_eval} ; on s'appuie ici sur sa conclusion : le
top-down borne l'ensemble des objets saisissables, et c'est lui qu'il faut
rendre robuste.

\subsection{Construction géométrique}

Pour un objet de centre \((x, y, z)\) dans le repère base, la classe
\texttt{TopDownGrasp} construit les trois poses en empilant la même rotation
\(R(\theta)\) (pince vers le bas, \textit{yaw} \(\theta\)) sur trois hauteurs :

\begin{align}
  T_{\text{approach}} &= T_{\text{base}}^{\text{gripper}}\!\big(g_x,\, g_y,\, z_{\text{grasp}} + h_{\text{approach}},\, R(\theta)\big) \\
  T_{\text{grasp}}    &= T_{\text{base}}^{\text{gripper}}\!\big(g_x,\, g_y,\, z_{\text{grasp}},\, R(\theta)\big) \\
  T_{\text{retract}}  &= T_{\text{base}}^{\text{gripper}}\!\big(g_x,\, g_y,\, z_{\text{grasp}} + h_{\text{retract}},\, R(\theta)\big)
\end{align}

avec \(h_{\text{approach}} = 8\) cm et \(h_{\text{retract}} = 10\) cm. Le couple
\((g_x, g_y)\) n'est pas exactement \((x, y)\) : il intègre un décalage latéral
(section~\ref{sub:plan_lateral}). La hauteur de prise \(z_{\text{grasp}}\) n'est
pas non plus le \(z\) détecté par la stéréo : elle est ancrée sur la table
(section~\ref{sub:plan_depth}). Les sous-sections suivantes établissent chacune
de ces composantes ; on commence par la rotation, qui porte l'essentiel de la
géométrie.

\subsection{Rotation \texorpdfstring{$R(\theta)$}{R(theta)} : pince verticale alignée}
\label{sub:plan_rotation}

On cherche une orientation \(R\) de l'effecteur telle que l'axe de fermeture des
doigts (noté \(Z_{\text{pince}}\)) pointe vers le bas, c'est-à-dire vers
\(-Z_{\text{base}}\), tout en laissant libre une orientation horizontale réglable
par un angle de \textit{yaw} \(\theta\) (l'angle de l'articulation
\texttt{wrist\_roll}). Posons la convention de la pince : \(Z_{\text{pince}}\)
est l'axe de fermeture, \(X_{\text{pince}}\) l'axe « vue de face », et
\(Y_{\text{pince}}\) complète le trièdre (règle de la main droite).

Imposer \(Z_{\text{pince}} = -Z_{\text{base}}\) ne suffit pas à fixer \(R\) : il
reste un degré de liberté, le \textit{yaw} autour de l'axe vertical. On le
paramètre en plaçant l'axe « vue de face » dans le plan horizontal,
\(X_{\text{pince}} = (\cos\theta,\, \sin\theta,\, 0)\). Le troisième vecteur
colonne se déduit par produit vectoriel,
\(Y_{\text{pince}} = Z_{\text{pince}} \times X_{\text{pince}} = (\sin\theta,\, -\cos\theta,\, 0)\).
En empilant ces trois vecteurs en colonnes :

\begin{equation}
  R(\theta) = \begin{pmatrix}
    \cos\theta &  \sin\theta &  0 \\
    \sin\theta & -\cos\theta &  0 \\
    0          &  0          & -1
  \end{pmatrix}
\end{equation}

Une rotation \emph{propre} doit vérifier \(R R^\top = I_3\) (orthonormalité) et
\(\det R = +1\) (sens direct conservé). L'orthonormalité tient parce que les trois
colonnes sont unitaires et deux à deux orthogonales par construction. Pour le
déterminant, en développant par la dernière ligne :
\[
  \det R = (-1)\,\big(\cos\theta \cdot (-\cos\theta) - \sin\theta \cdot \sin\theta\big)
         = (-1)\,(-\cos^2\theta - \sin^2\theta) = +1 .
\]
Géométriquement, \(R(\theta)\) est la composée d'une rotation de \(\pi\) autour de
\(X\) (qui retourne \(Z\) et \(Y\), d'où la pince qui regarde vers le bas) suivie
d'une rotation de \(\theta\) autour du nouvel axe vertical. Le point qui mérite
d'être souligné : inverser \emph{un seul} axe (par exemple poser
\(Z_{\text{pince}} = -Z_{\text{base}}\) sans corriger un second axe) produirait
une réflexion, de déterminant \(-1\), inutilisable comme rotation d'un solide.
C'est précisément la correction sur \(Y_{\text{pince}}\) (le signe \(-\) en
position centrale) qui ramène \(R\) dans le groupe des rotations propres
\(SO(3)\). Le premier des self-tests
(section~\ref{sec:plan_validation}) vérifie numériquement ces deux propriétés.

\subsection{Orientation \texttt{wrist\_roll} : machoires en travers du petit axe}
\label{sub:plan_yaw}

Reste à fixer le \textit{yaw} \(\theta\). La règle est unique, sans cas
particulier par objet : on aligne les machoires \textbf{en travers du petit axe}
de l'empreinte de l'objet, c'est-à-dire perpendiculairement à son grand axe.
L'intuition est mécanique : pour un objet allongé (stylo, paquet de mouchoirs),
fermer la pince le long du grand axe la ferait glisser ; fermer en travers serre
le petit côté et offre une prise stable.

\paragraph{Source préférée : le grand axe en repère base.} La perception fournit,
quand l'empreinte au sol a une orientation fiable (objet nettement allongé),
l'angle \texttt{yaw\_base\_rad} de son grand axe \emph{déjà exprimé dans le
repère base} (projection rayon-plan du contour sur le plan de l'objet, puis ACP
2D ; cf chapitre~\ref{ch:perception}). On pose alors directement
\(\theta = \texttt{yaw\_base\_rad}\). Cet angle étant calculé en repère base, il
est correct quelle que soit l'orientation des caméras, ce qui permet de saisir
des objets posés \emph{en biais}.

\paragraph{Empreinte ronde ou indéterminée : yaw libre.} Lorsque la perception
classe l'objet comme \emph{debout} ou \emph{compact}, l'empreinte vue du dessus
est quasi circulaire et le grand axe n'a pas de sens. Imposer un angle serait
au mieux inutile, au pire nuisible. On marque alors le \textit{yaw} comme
\emph{libre} : c'est l'IK qui choisira l'angle minimisant la rotation du poignet
depuis la pose courante (chapitre~\ref{ch:control_eval}). Pour un objet rond,
tout angle saisit aussi bien ; ne rien imposer évite une rotation de poignet
gratuite.

\paragraph{Repli : l'angle du contour image.} Si la géométrie 3D n'est pas
disponible (détecteur V1 sans estimation de pose), on retombe sur une
approximation : l'angle principal du contour 2D dans l'image. La fonction
\texttt{yaw\_from\_contour} le calcule à partir des moments centraux d'ordre~2.
Posons le contour comme un nuage de points \((x_i, y_i)\) recentré sur son
centroïde ; les moments centraux sont
\[
  \mu_{20} = \sum_i x_i^2, \qquad
  \mu_{02} = \sum_i y_i^2, \qquad
  \mu_{11} = \sum_i x_i y_i .
\]
Ils forment la matrice de covariance 2D du nuage,
\(\Sigma = \left(\begin{smallmatrix}\mu_{20} & \mu_{11}\\ \mu_{11} & \mu_{02}\end{smallmatrix}\right)\),
dont le vecteur propre dominant donne le grand axe. L'angle de cet axe propre
admet une forme close :
\begin{equation}
  \theta = \tfrac{1}{2} \, \operatorname{atan2}\!\big(2\,\mu_{11},\; \mu_{20} - \mu_{02}\big) .
\end{equation}
Cette expression est la diagonalisation analytique d'une matrice symétrique
\(2\times2\) ; elle équivaut à une analyse en composantes principales du contour.
Pour une forme quasi-symétrique (cube, disque), \(\mu_{20} \approx \mu_{02}\) et
\(\mu_{11} \approx 0\) : l'axe principal devient mal défini. Un test de
dégénérescence (les deux écarts relatifs sous 5\,\%) renvoie alors
\(\theta = 0\), un alignement par défaut plutôt qu'une valeur de bruit. Il faut
être clair sur la portée de ce repli : l'angle obtenu est dans le repère
\emph{image}, pas dans le repère base, et il ignore l'inclinaison de l'objet.
C'est une estimation grossière, conservée uniquement pour le cas où la
perception 3D fait défaut ; la source préférée reste l'angle en repère base.

\subsection{Ouverture de pince adaptative}
\label{sub:plan_opening}

Une fois l'orientation fixée, il faut décider de combien ouvrir la pince avant
de descendre. Ouvrir systématiquement au maximum gaspille de la course et,
surtout, rapproche les doigts ouverts des objets voisins. On calcule donc une
ouverture \emph{adaptée à la largeur de l'objet}. À partir de la boîte
englobante 3D, on prend la plus petite des deux dimensions horizontales
(\(w = \min(b_x, b_y)\), la pince se refermant dans le plan horizontal), on
ajoute une marge \(m\) de chaque côté, et on rapporte à l'ouverture maximale
réelle de la pince :
\begin{equation}
  \text{pct} = \frac{w + 2m}{O_{\max}} \times 100 ,
\end{equation}
avec \(m = 10\) mm et \(O_{\max} = 150\) mm. La marge a un rôle précis : absorber
l'erreur de visée (de l'ordre de quelques millimètres entre l'IK et la
calibration) pour que le doigt \emph{fixe} dégage l'objet à la descente au lieu
de le percuter. La valeur résultante est bornée dans \([20, 100]\,\%\) : assez
d'ouverture pour passer autour de l'objet, sans ouvrir à fond inutilement.

\paragraph{Découverte : l'ouverture maximale réelle.} La valeur de
\(O_{\max}\) n'a pas été correcte d'emblée. Le code la fixait initialement à
50~mm. Avec cette valeur, la formule saturait à 100\,\% dès qu'un objet dépassait
\(\sim\!26\) mm : pour le cube de 30~mm comme pour à peu près tout, la pince ne
savait que s'ouvrir au maximum, et l'ouverture « adaptative » ne différenciait
plus rien. La mesure terrain de la course réelle a donné 150~mm. Avec cette
valeur corrigée, un objet de 30~mm donne \((30 + 2\times10)/150 \approx 33\,\%\),
soit une cinquantaine de millimètres d'écartement, et l'ouverture distingue
enfin les objets selon leur taille.

\subsection{Décalage latéral vers le doigt fixe}
\label{sub:plan_lateral}

La pince du SO-101 est \textbf{asymétrique} : un doigt est fixe, l'autre se
referme contre lui. Centrer la prise sur le centroïde de l'objet a un effet
pervers : à la descente, le doigt fixe heurte l'objet \emph{avant} la fermeture,
le déplace, et la saisie rate. La parade est de décaler le centre de prise dans
la direction \emph{opposée} au doigt fixe, pour que l'objet finisse appuyé contre
ce doigt fixe et que le doigt mobile vienne simplement l'écraser. Si
\(\mathbf{d}\) est la direction du doigt fixe dans le repère pince et \(\ell\)
l'amplitude du décalage, le déplacement appliqué en repère base est
\[
  \Delta_{\text{base}} = R(\theta) \, (-\mathbf{d}) \, \ell ,
\]
appliqué identiquement à \textit{approach}, \textit{grasp} et \textit{retract}
pour garder l'alignement vertical. La valeur \emph{nominale} est \(\ell = 8\) mm,
utile pour les objets à faces plates (type cube) où un bord franc cale contre le
doigt fixe. En \emph{déploiement}, la valeur par défaut est ramenée à \(0\) :
pour un objet rond ou fin, le décalage ne sert à rien et risque au contraire de
faire manquer l'objet. On l'active explicitement (option
\texttt{--grasp-lateral-offset}) quand l'objet le justifie.

\subsection{Profondeur de prise ancrée sur la table}
\label{sub:plan_depth}

La hauteur de prise \(z_{\text{grasp}}\) mérite un traitement à part, parce que
l'axe \(Z\) est le plus bruité de la triangulation stéréo. Le
chapitre~\ref{ch:perception} l'a établi quantitativement : l'erreur de profondeur
croît comme \(Z^2\), et une erreur de quelques pixels en disparité se traduit par
plusieurs millimètres à un demi-mètre. Utiliser ce \(Z\) bruité comme hauteur de
prise expose à un risque concret : si l'estimation tombe sous le niveau réel, la
pince force dans la table.

On exploite plutôt un fait géométrique sûr : l'objet \emph{repose} sur la table,
et le niveau de la table est connu (\(z_{\text{table}} \approx 0\), la plaque du
robot étant posée dessus). On dérive donc la hauteur de prise du plan de la table
augmenté de la moitié de la hauteur estimée de l'objet (sa mi-hauteur) :
\begin{equation}
  z_{\text{grasp}} = z_{\text{table}} + \frac{h_{\text{objet}}}{2} ,
\end{equation}
et l'on n'utilise le \(Z\) détecté qu'en repli, faute de boîte englobante. Un
garde-fou borne le résultat : on ne descend jamais sous
\(z_{\text{table}} + 5\) mm. La profondeur de prise repose ainsi sur la dimension
\emph{verticale} de l'objet (relativement fiable, voir la triangulation du sommet
au chapitre~\ref{ch:perception}) et sur le plan de la table, jamais sur le
\(Z\) absolu de la stéréo.

\subsection{Découverte : la convention de fermeture à 90\textdegree}
\label{sub:plan_90deg}

Avec l'orientation, l'ouverture, le décalage et la profondeur en place, les
premières saisies réelles \emph{échouaient pourtant systématiquement} sur
l'alignement. Sur un objet allongé dont la perception donnait pourtant un grand
axe correct, la pince se présentait perpendiculaire à l'angle attendu et
manquait l'objet ou le poussait. L'hypothèse de départ — un \textit{yaw} mal
calculé en perception — ne tenait pas : les angles loggués étaient cohérents avec
ce que montraient les images.

Pour trancher, j'ai instrumenté la pipeline d'un \emph{diagnostic par caméra},
loggant l'angle de l'objet tel que vu par \texttt{cam\_0} et \texttt{cam\_1}
(\texttt{yaw\_cam0\_deg}, \texttt{yaw\_cam1\_deg}), puis comparé à l'angle de
fermeture effectivement réalisé par les machoires. L'écart était reproductible :
les doigts se fermaient à \(90\)\textdegree{} de la convention nominale du code.
Autrement dit, l'axe que le code croyait être l'axe de fermeture des machoires
était en réalité l'axe perpendiculaire, du fait du montage physique de la pince.
Un test empirique a confirmé l'hypothèse : en ajoutant un décalage de
\(90\)\textdegree{} au \textit{yaw} de prise (option \texttt{--grasp-yaw-offset 90}),
les saisies parallèles à \(X\) comme parallèles à \(Y\) ont réussi du premier
coup, avec un couple de maintien dépassant 300. La valeur \(0\) échouait à chaque
fois. Le décalage de \(90\)\textdegree{} est donc adopté comme valeur de
déploiement, tandis que la convention nominale (\(0\)) reste celle des self-tests.
Cet épisode illustre une difficulté propre au travail sur matériel réel : une
convention de repère parfaitement cohérente \emph{en code} peut être décalée par
rapport au \emph{montage physique}, et seul un essai instrumenté le révèle.

\subsection{Filtres et garde-fous}

Le \texttt{TopDownGrasp} rejette les objets trop hauts pour une saisie verticale,
quand le sommet estimé de l'objet dépasse \texttt{max\_object\_height\_m}
(\(0{,}12\) m par défaut) : au-delà, la pince entrerait en collision avec l'objet
avant d'atteindre sa mi-hauteur. Dans le brouillon initial, ce rejet motivait une
extension « \textit{side grasp} » pour les objets hauts ; la
section~\ref{sec:plan_horsportee} explique pourquoi cette piste a été remplacée
par un résultat de discussion, et non par du travail à venir.

\section{Au-delà du top-down : la prise latérale}
\label{sec:plan_horsportee}

Le brouillon de ce mémoire présentait la prise latérale (\textit{side grasp})
comme une extension à implémenter pour saisir les objets trop hauts pour une
attaque verticale. Cette présentation n'est plus exacte, et il serait malhonnête
de la maintenir : la prise latérale a été \emph{étudiée}, et le résultat est
qu'elle est mal conditionnée sur le SO-101, pas qu'elle reste à coder.

Le raisonnement complet est mené au chapitre~\ref{ch:control_eval}. En résumé :
les trois articulations de tangage du bras plient toutes dans un même plan
vertical, si bien que l'effecteur ne peut pas pointer son approche \emph{hors}
de ce plan. Une fermeture horizontale « vraie » autour du flanc vertical d'un
objet demande précisément ce type d'orientation. L'évaluation numérique sur un
grand nombre de configurations le confirme : l'IK ne réalise une telle fermeture
qu'avec plusieurs centimètres d'erreur, contre une erreur sous le millimètre pour
un top-down accessible. Le résultat n'est pas que la prise latérale soit
catégoriquement impossible — un bras à 6~degrés de liberté lèverait la
contrainte — mais qu'elle est \emph{mal conditionnée} et peu fiable sur 5~degrés.

La conséquence pour la planification est directe et, plutôt qu'une lacune, c'est
un cadrage : le top-down s'impose comme l'approche robuste, et c'est lui qui
\emph{borne} l'ensemble des objets saisissables (ceux accessibles par le dessus,
sous la hauteur limite). On renvoie donc le lecteur au
chapitre~\ref{ch:control_eval} pour la démonstration chiffrée, sans présenter de
\textit{side grasp} comme un développement futur.

\section{Interface avec le module de contrôle}
\label{sec:plan_interface}

La \texttt{GraspPose} produite ici est l'unique contrat entre la planification et
le contrôle : trois poses \(SE(3)\) et deux paramètres d'ouverture. Tout ce qui
précède (orientation, ouverture, décalage, profondeur, convention de fermeture)
est résolu \emph{avant} que le contrôle n'intervienne. Le module \textit{control}
(chapitre~\ref{ch:control_eval}) consomme ce contrat et réalise :
\begin{enumerate}
  \item pour chaque pose \(4\times4\) de \(SE(3)\), l'inversion de la cinématique
        (IK) pour obtenir les 5~angles articulaires ;
  \item l'interpolation des trajectoires entre poses successives ;
  \item l'envoi des commandes moteur via le bus Feetech ;
  \item la gestion de l'ouverture et de la fermeture de la pince.
\end{enumerate}

Un détail de l'interface mérite d'être noté, car il fait le lien avec le
\textit{yaw} libre de la section~\ref{sub:plan_yaw} : quand la perception fournit
en cours d'approche une mesure d'orientation plus fiable (typiquement
\texttt{cam\_2}, vue rapprochée et quasi verticale), la fonction
\texttt{reorient\_grasp\_pose} réoriente la \texttt{GraspPose} sur ce nouvel
angle en préservant la position de l'objet, le décalage latéral et la correction
de convention. C'est l'amorce du raffinement en boucle fermée que détaille le
chapitre suivant.

\section{Validation et tests}
\label{sec:plan_validation}

Le module \texttt{src/planning/grasp.py} embarque 6~self-tests
(\texttt{python -m src.planning.grasp}), chacun ciblant une des briques
géométriques établies plus haut :
\begin{enumerate}
  \item \textbf{Rotation propre} : \(R(\theta)\) orthonormale, \(\det R = +1\),
        et \(Z_{\text{pince}}\) pointant bien vers \(-Z_{\text{base}}\)
        (section~\ref{sub:plan_rotation}).
  \item \texttt{yaw\_from\_contour} : rectangle horizontal \(\rightarrow \theta = 0\),
        rectangle vertical \(\rightarrow \theta = \pm\pi/2\)
        (section~\ref{sub:plan_yaw}).
  \item \texttt{TopDownGrasp.plan} (décalage nul) : positions \((x, y, z)\)
        correctement propagées sur les 3 poses, pince vers le bas.
  \item \texttt{TopDownGrasp.plan} (décalage latéral 8~mm et ancrage table) :
        décalage en \(Y\) vers le doigt fixe et hauteur de prise à
        \(z_{\text{table}} + h/2\) (sections~\ref{sub:plan_lateral}
        et~\ref{sub:plan_depth}).
  \item \textbf{Filtre objet trop haut} : un objet de 15~cm retourne
        \texttt{None}.
  \item \textbf{Classe abstraite} \texttt{GraspStrategy} : \texttt{TypeError} si
        instanciation directe (garantit que toute stratégie future implémente
        l'interface).
\end{enumerate}

Les six tests passent. Ils valident la géométrie de manière isolée, hors robot ;
la validation \emph{en boucle fermée}, sur le matériel, relève du chapitre
suivant, où une campagne préliminaire informelle est rapportée et où la campagne
formelle reste à mener.

\todoF{Ajouter une figure illustrant les 3 poses (approach / grasp / retract) au-
dessus d'un objet schématique, avec le repère pince ($Z_{\text{pince}}$ vers le
bas, $X_{\text{pince}}$ horizontal) et le décalage latéral vers le doigt fixe.}

\todoF{Voix de Maxence : un mot court sur l'effet « eurêka » du diagnostic 90\textdegree{}
(section~\ref{sub:plan_90deg}) — combien de saisies ratées avant de penser à
instrumenter par caméra ? C'est un bon exemple du ton découverte demandé.}
